% ============================================================
\chapter*{Preface}
\addcontentsline{toc}{chapter}{Preface}
% ============================================================

\section*{Objectives and motivation}

Measure theory and integration form one of the fundamental pillars of modern
mathematical analysis. Born from the pioneering work of \'Emile Borel and Henri
Lebesgue at the beginning of the 20th century, this theory profoundly transformed
our understanding of the concepts of length, area, volume, and more generally the
``size'' of a set. The Lebesgue integral, which arises from it, considerably extends
the Riemann integral and provides the natural framework for probability theory,
functional analysis, partial differential equations, and many other branches of
mathematics.

This course is aimed at graduate students (Master's level) in pure or applied
mathematics. It assumes a solid command of real analysis (sequences, series,
continuity, differentiability, Riemann integral) as well as basic notions of general
topology. Familiarity with linear algebra and normed vector spaces is also desirable.

\section*{Historical context}

The Riemann integral, developed in the 1850s, has well-known limitations. Consider
the Dirichlet function:
\[
f(x) = \mathbf{1}_{\Q}(x) = \begin{cases} 1 & \text{if } x \in \Q, \\ 0 & \text{if } x \notin \Q. \end{cases}
\]
This function is not Riemann-integrable, yet intuitively, since $\Q$ is
``negligible'' in $\R$, one would expect $\int_0^1 f = 0$.

Similarly, if $(f_n)$ is a sequence of Riemann-integrable functions converging
pointwise to a function $f$, it is generally not true that
\[
\lim_{n \to \infty} \int_a^b f_n(x)\, dx = \int_a^b f(x)\, dx,
\]
even when the right-hand side makes sense. For the interchange of limit and integral
to be valid in the Riemann sense, one needs \emph{uniform} convergence, which is a
very restrictive condition.

The Lebesgue integral solves both problems elegantly: the Dirichlet function becomes
integrable (with integral zero), and the dominated convergence theorem allows the
interchange of limit and integral under much weaker hypotheses than uniform
convergence.

\section*{Guiding ideas}

Lebesgue's approach rests on a fundamentally different idea from Riemann's. While
Riemann partitions the \emph{domain} of the function into small intervals, Lebesgue
partitions the \emph{range} of the function and measures the preimages:

\begin{intuition}[title=\textbf{Lebesgue's philosophy}]
Imagine you need to count a large sum of money composed of coins of different values.
Riemann's method takes the coins one by one, in the order they are found. Lebesgue's
method first sorts the coins by value, then counts how many there are of each type.
The result is the same, but the second method is more flexible and applies to more
situations.
\end{intuition}

For this approach to work, one must be able to ``measure'' the preimages
$f^{-1}([a,b))$. This requires:
\begin{enumerate}
    \item A theory of \textbf{measurable sets} ($\sigma$-algebras), which determines
    which sets can be measured.
    \item A theory of \textbf{measures}, which assigns a ``size'' to each of these sets.
    \item A definition of \textbf{measurable functions}, whose preimages of
    reasonable sets are measurable.
    \item The construction of the \textbf{integral} itself, by approximation via
    step functions (\emph{simple functions}).
\end{enumerate}

\section*{Course outline}

The course is organized into twelve chapters following this logical progression:

\begin{description}
    \item[Chapter 1.] $\sigma$-Algebras and measurable spaces: the fundamental
    structures. Generated $\sigma$-algebras, Borel $\sigma$-algebra.
    \item[Chapter 2.] Measures: definitions, examples, $\sigma$-additivity,
    finite and $\sigma$-finite measures, continuity lemma.
    \item[Chapter 3.] Carath\'eodory's theorem: construction of measures via outer
    measures, extension from a pre-measure to a complete measure.
    \item[Chapter 4.] Lebesgue measure on $\R^n$: construction, invariance
    properties, non-measurable sets.
    \item[Chapter 5.] Measurable functions: definitions, stability under algebraic
    operations and limits, approximation by simple functions.
    \item[Chapter 6.] Lebesgue integral: three-step construction (simple functions,
    non-negative functions, integrable functions).
    \item[Chapter 7.] Convergence theorems: monotone convergence (Beppo Levi),
    Fatou's lemma, dominated convergence (Lebesgue).
    \item[Chapter 8.] $L^p$ spaces: H\"older and Minkowski inequalities,
    completeness, separability, duality.
    \item[Chapter 9.] Signed measures: Hahn and Jordan decompositions, total
    variation.
    \item[Chapter 10.] Radon--Nikodym theorem: absolute continuity, Radon--Nikodym
    derivative, Lebesgue decomposition.
    \item[Chapter 11.] Product measures: product $\sigma$-algebra, Fubini and
    Tonelli theorems.
    \item[Chapter 12.] Radon measures and the Riesz representation theorem.
\end{description}

\section*{Conventions and notation}

\begin{itemize}
    \item $\N = \{0, 1, 2, \ldots\}$ denotes the set of natural numbers (including
    zero).
    \item $\R^+ = [0, +\infty)$ and $\overline{\R}^+ = [0, +\infty]$.
    \item The extended real line is $\overline{\R} = [-\infty, +\infty]$.
    \item $\mathbf{1}_A$ denotes the indicator (characteristic) function of $A$.
    \item $f^+ = \max(f, 0)$, $f^- = \max(-f, 0)$, $f = f^+ - f^-$,
    $\abs{f} = f^+ + f^-$.
    \item $A_n \uparrow A$ means $(A_n)$ is increasing and $A = \bigcup_n A_n$.
    \item $A_n \downarrow A$ means $(A_n)$ is decreasing and $A = \bigcap_n A_n$.
    \item ``a.e.'' means ``almost everywhere'' (outside a set of measure zero).
    \item $(X, \mathcal{A}, \mu)$ denotes a measure space.
    \item $\mathcal{B}(\R^n)$ denotes the Borel $\sigma$-algebra on $\R^n$.
    \item $\lambda^n$ denotes the Lebesgue measure on $\R^n$ ($\lambda = \lambda^1$).
\end{itemize}

\section*{Prerequisites}

\begin{itemize}
    \item \textbf{Real analysis}: properties of $\R$, numerical sequences and series,
    continuity, differentiability, Riemann integral, series of functions, uniform
    convergence.
    \item \textbf{General topology}: metric spaces, open and closed sets, compactness,
    connectedness.
    \item \textbf{Linear algebra}: vector spaces, norms, Banach spaces (basic notions).
    \item \textbf{Set theory}: set operations, axiom of choice (for the construction
    of non-measurable sets).
\end{itemize}

\section*{References}

\begin{itemize}
    \item H.~L.~Royden and P.~M.~Fitzpatrick, \emph{Real Analysis}, 4th ed.,
    Pearson, 2010.
    \item W.~Rudin, \emph{Real and Complex Analysis}, 3rd ed., McGraw-Hill, 1987.
    \item G.~B.~Folland, \emph{Real Analysis: Modern Techniques and Their
    Applications}, 2nd ed., Wiley, 1999.
    \item P.~Halmos, \emph{Measure Theory}, Springer, 1974.
    \item R.~M.~Dudley, \emph{Real Analysis and Probability}, Cambridge University
    Press, 2002.
    \item E.~M.~Stein and R.~Shakarchi, \emph{Real Analysis: Measure Theory,
    Integration, and Hilbert Spaces}, Princeton University Press, 2005.
    \item R.~L.~Schilling, \emph{Measures, Integrals and Martingales}, Cambridge
    University Press, 2005.
    \item A.~N.~Kolmogorov and S.~V.~Fomin, \emph{Introductory Real Analysis},
    Dover, 1975.
\end{itemize}

\bigskip
\hfill \emph{The Author}

\hfill \emph{March 2026}
